\documentclass[12pt]{article}
\usepackage[margin=1.25in]{geometry}
\usepackage{setspace}
\doublespacing
\frenchspacing
\usepackage{amsmath,amssymb}
\usepackage{booktabs}
\usepackage{longtable}
\usepackage{graphicx}
\usepackage{float}
\usepackage{enumitem}
\usepackage[dvipsnames]{xcolor}
\usepackage[hidelinks]{hyperref}
\hypersetup{
  colorlinks=false,
  pdfborder={0 0 0},
  allbordercolors={1 1 1}
}
\usepackage{natbib}
\bibpunct{(}{)}{;}{a}{,}{ }
\usepackage{microtype}
\usepackage{caption}
\captionsetup[table]{
  labelfont={bf},
  textfont={normalsize},
  labelsep=period,
  justification=raggedright,
  singlelinecheck=false,
  skip=6pt
}
\captionsetup[figure]{font=small,labelfont=bf,labelsep=period}
\renewcommand{\tablename}{Table}
\renewcommand{\figurename}{Figure}
\usepackage{titlesec}

%% ---- Section/subsection formatting matching the main paper -----------
\titleformat{\section}
  {\normalfont\large\bfseries\color{black}}
  {\thesection.}{0.5em}{}
\titlespacing*{\section}{0pt}{18pt}{8pt}

\titleformat{\subsection}
  {\normalfont\normalsize\bfseries\color{black}}
  {\thesubsection}{0.5em}{}
\titlespacing*{\subsection}{0pt}{12pt}{4pt}

\titleformat{\subsubsection}
  {\normalfont\normalsize\itshape\color{black}}
  {}{0em}{}
\titlespacing*{\subsubsection}{0pt}{8pt}{3pt}

\usepackage{xr}
\externaldocument{main}
\renewcommand{\refname}{References}

\hyphenation{confidence taxonomy}

\color{black}

\newcommand{\JF}{\textit{Journal of Finance}}
\newcommand{\RFS}{\textit{Review of Financial Studies}}
\newcommand{\JFE}{\textit{Journal of Financial Economics}}

\title{Online Appendix to ``Do Finance Journals Pick Sides? Evidence from Government Intervention Research''}
\author{}
\date{}

\begin{document}
\maketitle

\section*{Overview}

This Online Appendix collects three supporting materials referenced in the main paper:
the full keyword taxonomy used at the first-pass classification stage
(Online Appendix~OA.1), the verbatim LLM prompts used to assign in-scope
and direction codes (Online Appendix~OA.2), and the sample-frame
precision and recall audit (Online Appendix~OA.3).  Section, table, and
figure numbers in this Online Appendix begin with ``OA''.  References to
sections of the main paper retain their main-paper numbering.

\renewcommand{\thesection}{OA.\arabic{section}}
\setcounter{section}{0}
\renewcommand{\thetable}{OA.\arabic{table}}
\setcounter{table}{0}

\section{Keyword Taxonomy}

Table~\ref{tab:keywords} lists all terms used in Pass~1 of the classification
pipeline.

\begin{table}[H]
\centering
\caption{In-Scope Classification: Keyword \mbox{Taxonomy}}
\label{tab:keywords}
\small
\renewcommand{\arraystretch}{1.15}
\begin{tabular}{@{}p{3.8cm}p{10cm}@{}}
\toprule
\textbf{Category} & \textbf{Keywords} \\
\midrule
Financial regulation &
  regulation, regulatory, Dodd-Frank, Basel, Glass-Steagall,
  capital requirements, leverage ratio, bank capital \\
\addlinespace[3pt]
Securities law &
  SEC, securities and exchange, disclosure requirement, Reg FD,
  Sarbanes-Oxley, SOX, insider trading law, short-selling ban,
  short sale ban, short-sale restriction \\
\addlinespace[3pt]
Banking policy &
  FDIC, deposit insurance, lender of last resort, bailout, TARP,
  stress test, bank resolution, government guarantee \\
\addlinespace[3pt]
Monetary/fiscal policy &
  quantitative easing, QE, Fed intervention, government subsidy,
  public guarantee, state guarantee, fiscal stimulus \\
\addlinespace[3pt]
Market microstructure rules &
  circuit breaker, tick size, trading halt, uptick rule,
  limit-down, limit-up, trading curb \\
\addlinespace[3pt]
Corporate governance law &
  board independence requirement, say-on-pay, proxy access,
  mandatory audit, governance mandate, clawback provision,
  mandatory disclosure \\
\addlinespace[3pt]
General intervention signals &
  government intervention, government policy, public policy,
  law and finance, legal protection, shareholder rights,
  creditor rights, investor protection, financial regulation reform \\
\bottomrule
\end{tabular}
\end{table}

\noindent Ambiguous standalone terms (\textit{regulation}, \textit{regulatory},
\textit{SEC}) require co-occurrence with at least one additional in-scope
keyword before the paper is flagged as a candidate.

%% ============================================================


\section{LLM Prompts}

\subsection{In-Scope Classification Prompt}

The prompt receives the paper's title and abstract.
System instruction: ``You are a research assistant helping to classify
finance research papers. You must return valid JSON only, no other text.''
User message:

\begin{quote}
\small
A finance research paper is `in scope' for a study of directional bias
in government intervention research if and only if its PRIMARY research
question is the CAUSAL EFFECT of a specific government law, regulation,
or public policy on financial markets, firms, or investors.

Exclusion criteria (mark in\_scope=false):
\newline, Purely theoretical papers (no empirical test of a policy)
\newline, Papers that only DESCRIBE a regulation without estimating its effect
\newline, Papers whose main topic is private firm/market behaviour, with policy as incidental context
\newline, Papers studying central bank communication WITHOUT a specific policy instrument
\newline, Papers whose primary policy variable is a rule imposed by a stock exchange
(NYSE, NASDAQ, etc.) without a specific government or regulatory mandate
(e.g., exchange-designed trading halts, voluntary tick size conventions,
exchange fee structures)

Return JSON with exactly these keys: in\_scope (true or false),
confidence (``high'' $|$ ``medium'' $|$ ``low''), reason (one sentence).
\end{quote}

\subsection{Direction Coding Prompt}

The deployed pipeline has two paths, selected per paper by the
direction-coding classifier.  When the abstract is present and
contains at least 20 words, the model receives the
\textit{abstract-only} prompt below; the title is not
concatenated with the abstract or otherwise passed.  When the abstract is missing or
shorter than 20 words, the model instead receives the \textit{title-only}
fallback prompt also shown below.  All 291 papers in the final analysis
sample have abstracts of at least 20 words, so the title-only fallback
was not invoked for any analysis-sample paper; the validation samples
also use the same two-path logic.

System instruction (both prompts): ``You are a research assistant
classifying finance paper findings. Return valid JSON only, no other
text.''

\paragraph{Abstract-only prompt (used for all 291 analysis-sample papers).}
\begin{quote}
\small
The following is the abstract of a finance paper that studies the effect
of a government law, regulation, or public policy.

Classify the MAIN empirical finding:
\newline\quad +1 \enspace The government intervention produced a POSITIVE / BENEFICIAL
outcome (e.g.\ improved market quality, reduced systemic risk, better investor
protection, increased firm value, lower cost of capital)
\newline\quad -1 \enspace The government intervention produced a NEGATIVE / HARMFUL
outcome (e.g.\ reduced liquidity, higher costs, unintended consequences, welfare losses)
\newline\quad\phantom{+}0 \enspace NULL, MIXED, or INCONCLUSIVE finding (no statistically
significant effect, evidence on both sides, or the paper explicitly refuses to
take a directional stance)

Coding rules:
\newline, Code the PRIMARY intervention and PRIMARY outcome (named in title or first sentence)
\newline, If paper is purely theoretical with no empirical test,
\newline\phantom{,}\quad set direction=0, confidence=low
\newline, Do NOT infer direction from normative framing; code the sign of the estimated coefficient

Abstract:\\
\texttt{[ABSTRACT TEXT HERE]}

Return JSON with exactly: direction (1 $|$ -1 $|$ 0),
confidence (``high'' $|$ ``medium'' $|$ ``low''), rationale (one sentence).
\end{quote}

The instruction ``named in title or first sentence'' is preserved from
the deployed prompt verbatim; because only the abstract is passed in
this path, the model interprets ``title'' as title-like leading material
that some abstracts include in their first sentence (e.g., naming the
specific statute being studied).  Section~\ref{subsec:direction} of the
main paper paraphrases the spirit of the rule, coding the sign of the estimated
effect relative to the stated regulatory goal, which the prompt conveys
through the example outcomes for $+1$ (regulatory objectives achieved:
improved market quality, reduced systemic risk, etc.) and $-1$
(regulatory side-effects: reduced liquidity, higher costs, etc.).

\paragraph{Title-only fallback prompt (documented for pipeline completeness; not invoked for any paper in the 291-paper analysis sample).}
\begin{quote}
\small
The following is the title of a finance paper that studies the effect of
a government law, regulation, or public policy. No abstract is available.

Based on the title alone, infer the likely direction of the main empirical finding:
\newline\quad +1 \enspace The intervention likely produced a POSITIVE / BENEFICIAL outcome
\newline\quad -1 \enspace The intervention likely produced a NEGATIVE / HARMFUL outcome
\newline\quad\phantom{+}0 \enspace NULL, MIXED, or cannot be determined from the title alone

Coding rules:
\newline, Default to 0 (null) if the title is neutral or provides no directional signal
\newline, Only assign $+1$ or $-1$ if the title clearly implies a directional finding
\newline, Set confidence to ``low'' unless the directional signal is unambiguous

Title:\\
\texttt{[TITLE TEXT HERE]}

Return JSON with exactly: direction (1 $|$ -1 $|$ 0),
confidence (``high'' $|$ ``medium'' $|$ ``low''), rationale (one sentence).
\end{quote}

%% ============================================================


\section{Sample Frame Validation: Precision and Recall}
%% ============================================================

The selection pipeline narrows 39,136 documents to 1,849 keyword
candidates and then to 291 final papers (Table~\ref{tab:sample_selection} of the main paper).
Differential exclusion at any stage, between NBER and journal papers, or
between null and directional papers, could bias the within-direction
corpus-membership shares that our main inference relies on.  This appendix
reports three precision/recall estimates that bound this concern.

\subsection*{LLM precision on its in-scope decisions}

We re-evaluated 68 papers that the LLM flagged as in-scope to test whether
the in-scope label was correct under the boundary rule of
Section~\ref{subsec:inscope} of the main paper.  The author overrode 10 of the 68, giving an
implied LLM precision (against the manual gold standard) of
$58/68 = 85.3\%$.  Across the full keyword-positive pool, the LLM flagged
217 papers as in-scope at the first pass; manual review trimmed this to
185 ($32$ of the initial $217$ were overridden: 10 NBER and 22 published,
with the published-side split JF: 18, RFS: 2, JFE: 2).  The NBER refilter
then promoted 126 additional papers to in-scope, of which 20 were overridden
on second-pass review and 106 survived, yielding the final 291-paper sample
($217 - 32 + 106 = 291$, equivalently $217 + 126 - 52 = 291$ with the
combined override count).  Every per-paper override is recorded in the
replication package; the boundary criterion (policy as primary research
question vs.\ as identification device for a non-policy outcome) is
operationalized in Section~\ref{subsec:inscope} of the main paper.

\subsection*{LLM recall on the keyword-positive pool}

The NBER refilter re-coded 803 NBER papers that the LLM had flagged as
low-confidence out-of-scope; 126 were promoted to in-scope and 106
survived the same manual review used elsewhere.  Treating
manual + refilter as the gold standard restricted to the
NBER subset, the LLM's first-pass recall on this low-confidence pool is
$1 - 126/803 = 84.3\%$, comparable to its precision.  The published-side
LLM decisions were not subjected to a symmetric refilter because the
journal denominators (8,924 articles) are small enough that the keyword
filter retains all plausible candidates at the first stage, but the
symmetry between the precision and recall point estimates suggests no
strong directional asymmetry in LLM errors within the keyword-positive
pool.

\subsection*{Keyword-filter recall: audit of keyword-negative papers}

Stage~1 (the keyword filter) is by design over-inclusive, but its recall
on truly in-scope papers cannot be inferred from
Table~\ref{tab:sample_selection} of the main paper alone.  We audit the keyword-negative
pool by sampling 200 papers (random seed 42) from the 37,287 documents
that fail the keyword filter and have abstracts of at least 20 words, and
re-running the in-scope LLM prompt on each.  Of the 200 sampled papers,
the LLM flagged $4$ as in-scope: false-negative rate
$2.0\%$ (95\% binomial CI: $0.5\%$ to $4.0\%$).  Extrapolated to the full
37,287-paper keyword-negative pool with usable abstracts, this implies
roughly $746$ in-scope papers (95\% CI: $186$ to $1{,}491$) that the
keyword filter misses; adjusting for the $\approx 85\%$ LLM precision
documented above brings the implied count of manually-validated
in-scope misses to roughly $634$ (95\% CI: $158$ to $1{,}267$).

The 4 audited keyword-misses concern papers studying specific government
policies (a tax rule change in commercial mortgage markets, voluntary
restraint arrangements with Japan, bank nationalization, foreign royalty
tax rates on R\&D) whose abstracts happen not to contain any term in our
seven-category taxonomy.  Per-source false-negative rates in the small
audit sample are not significantly differentiated
(JF: $1/14$, RFS: $1/7$, JFE: $0/15$, NBER: $2/164$), but the journal
side carries a slightly higher point estimate.  If a true differential
exists, the published-side denominator is modestly understated relative
to the NBER denominator; the implied bias on the within-direction
corpus-membership shares depends additionally on how the missed papers
distribute across direction categories, which we cannot estimate from
this small sample.

\subsection*{Interpretation}

Our main inference rests on the equality of direction shares across the
NBER and published samples within the selected corpus.  A
direction-neutral keyword miss leaves these shares unchanged; a
direction-correlated keyword miss would bias them.  The 200-paper
audit is too small to detect such differential exclusion if it exists.

%% ============================================================


\end{document}
